\documentclass[10pt,a4paper]{amsart}
\usepackage[margin=19mm]{geometry}
\usepackage{amsmath,amssymb,mathtools}
\usepackage[scr=rsfs,cal=euler]{mathalfa}
\usepackage{xfrac}
\usepackage[unicode,hidelinks,bookmarks=false]{hyperref}
\newcommand{\ie}{i.e.\ }
\newcommand{\cf}{cf.\ }
\newcommand{\resp}{respectively\ }
\newcommand{\Subsection}{\subsection}
\providecommand{\nobreakdash}{\nobreak\hbox{-}}
\setlength{\emergencystretch}{2em}
\multlinegap0pt
\usepackage{algorithm}\usepackage{algpseudocode}
\usepackage{xcolor}
\newtheorem{lemma}{Lemma}
\title{A Penalty Method for Non-Self-Adjoint Topology Optimization}
\author{Wei Gong, Yuanda Ye}
\date{Source: arXiv:2603.00962v1 (CC BY 4.0)}
\begin{document}
\maketitle

\noindent\emph{Source and licence.} A Penalty Method for Non-Self-Adjoint Topology Optimization. Version arXiv:2603.00962v1 (CC BY 4.0), distributed by arXiv under the Creative Commons Attribution 4.0 International licence, http://creativecommons.org/licenses/by/4.0/, as recorded in the arXiv licence metadata for this exact version and shown on the abstract page https://arxiv.org/abs/2603.00962v1. Full licence notice: \texttt{licenses/arxiv-2603.00962-CC-BY-4.0.txt}.\par\smallskip

\noindent\textit{Editorial scope.} Counted unit: the article's lemma F-second-derivative (source0.tex lines 370--374). The article's complete original proof of it is delivered with it (source0.tex lines 375--412, one \texttt{proof} environment of 38 lines). No mathematics is written, repaired or re-derived: the statement and the proof are the article's own lines, in the article's own notation, and no retained line is substituted or re-flowed.  Retained uncounted as the source-local context the proof needs: lines 174--193; lines 312--332.  Nothing was omitted from any retained span, so the package carries no omission record.  No retained line contains a citation, so the wrapper carries no bibliography and no bibliography entry.  No retained line contains a figure environment, an image-inclusion command or drawing code, so no image is delivered and the package declares no asset dependency.  Editorial formatting only, in addition: an AMS \texttt{amsart} preamble that restates the article's own declarations and macros used by the retained lines, namely the environments \texttt{lemma} on separate counters, together with the standard packages those lines need.  The retained environments print in this document as Lemma 1 in order of appearance, whereas the article prints the same items with the numbering of its own full document.  The complete native source of the article as distributed in the arXiv e-print for this exact version is bundled unchanged as \texttt{sources/195530/source0.tex}.
\begin{lemma}\label{stress_equ}
    The first-order optimality condition for the minimization problem 
    \begin{equation}\label{stress_problem}
        \min_{\sigma\in S_1}\int_{\Omega}E^{-1}(\chi)\sigma:\sigma\ dx
    \end{equation}
    can be characterized as follows
    \begin{equation}
    \left\{
    \begin{aligned}
    &\sigma = E(\chi)\varepsilon(\lambda) \quad &\text{in } \Omega, \\
    &\nabla \cdot \sigma = 0 \quad &\text{in } \Omega, \\
    &\lambda = 0 \quad &\text{on } \Gamma_D, \\
    %&\lambda = -\mu \quad &\text{on } \Gamma_N, \\
    &\lambda = -\mu,\quad \sigma \cdot \vec{n} = q_{\rm in} \quad &\text{on } \Gamma_N,
    \end{aligned}
    \right.
    \end{equation}
    where
    $$ S_1:=\{\sigma\in L^2(\Omega)^{d\times d}:\sigma_{i,j}=\sigma_{j,i},\ \nabla\cdot\sigma=0,\ \sigma\cdot\vec{n}|_{\Gamma_N}=q_{\rm in}\}.$$
\end{lemma}

\begin{lemma}\label{F-derivative}
    Define $$G(\chi) = \min_{(\hat\sigma,\hat\varrho)\in S} \tilde{g}(\chi,\hat\sigma,\hat\varrho):=\int_{\Omega}A(\chi)E_0^{-1}(\hat\sigma:\hat\sigma+\hat\varrho:\hat\varrho)\ dx,$$
    then from Lemma \ref{stress_equ} we have
    \begin{eqnarray}
        G(\chi) =\int_\Omega A(\chi)E_0^{-1}\big(\sigma(\chi):\sigma(\chi)+\varrho(\chi):\varrho(\chi)\big)\ dx.\nonumber
    \end{eqnarray}
    Here we denote the stress fields corresponding to \(\chi\) as \(\sigma(\chi):=\dfrac{1}{A(\chi)}{E}_0\varepsilon(u)\) and \(\varrho(\chi):=\dfrac{1}{A(\chi)}{E}_0\varepsilon(v)\) to emphasize their dependence on \(\chi\),
%    \begin{eqnarray}
%    (\sigma(\chi),\varrho(\chi)):=(\dfrac{1}{A(\chi)}{E}_0\varepsilon(u),\dfrac{1}{A(\chi)}{E}_0\varepsilon(v)),\nonumber    
%    \end{eqnarray}
    where $ u\in H_D^1(\Omega)^d $ and $ v\in H_D^1(\Omega)^d $ satisfy
    \begin{eqnarray} \label{u_weak}
    \int_{\Omega}\dfrac{1}{A(\chi)}{E}_0\varepsilon(u):\varepsilon(\hat{u})\ dx=\int_{\Gamma_N}q_{\rm in}\cdot\hat{u}\ ds\quad\forall\hat{u}\in H_D^1(\Omega)^d
    \end{eqnarray}
    and 
    \begin{eqnarray} \label{v_weak}
    \int_{\Omega}\dfrac{1}{A(\chi)}{E}_0\varepsilon(v):\varepsilon(\hat{v})\ dx=\int_{\Gamma_N}q_{\rm out}\cdot\hat{v}\ ds\quad\forall\hat{v}\in H_D^1(\Omega)^d.
    \end{eqnarray}
    Moreover, there holds
    $$\dfrac{\partial G}{\partial \chi}(\chi)(\hat{\chi})=\int_{\Omega}A'(\chi)(\hat{\chi}){E}_0^{-1}(\sigma(\chi):\sigma(\chi)+\varrho(\chi):\varrho(\chi))\ dx.$$
\end{lemma}

\begin{lemma}\label{F-second-derivative}
    The functional 
    $$g_1(\chi)=\min_{\sigma\in S_1}\int_{\Omega}E^{-1}(\chi)\sigma:\sigma\ dx$$ 
    is concave with respect to $ \chi $.
\end{lemma}

\begin{proof}
     It follows from the proof of Lemma \ref{F-derivative} that 
    \begin{align*}
g_1(\chi)=\int_{\Omega}\dfrac{1}{A(\chi)}{E}_0\varepsilon(u):\varepsilon(u)\ dx,
    \end{align*}
    where $u$ solves the state equation \eqref{u_weak}. Here, we do not employ the Lagrangian approach, but instead use the implicit function theorem. Denoting \( u' = \dfrac{\partial u}{\partial\chi}(\hat{\chi}) \) and \( u'' = \dfrac{\partial^2 u}{\partial\chi^2}(\hat{\chi}^2) \), by the chain rule, we can compute
    \begin{align*}
        \dfrac{\partial^2 g_1}{\partial\chi^{2}}(\hat{\chi}^2)=& 2\int_{\Omega}\dfrac{[A'(\chi)(\hat{\chi})]^2}{A(\chi)^{3}}{E}_0\varepsilon(u):\varepsilon(u)\ dx-4\int_{\Omega}\dfrac{A'(\chi)(\hat{\chi})}{A(\chi)^2}{E}_0\varepsilon(u):\varepsilon(u')\ dx\\
        &+2\int_{\Omega}\dfrac{1}{A(\chi)}{E}_0\varepsilon(u'):\varepsilon(u')\ dx+2\int_{\Omega}\dfrac{1}{A(\chi)}{E}_0\varepsilon(u):\varepsilon(u'')\ dx.
    \end{align*}
    Differentiating with respect to $\chi$ in \eqref{u_weak}, we obtain 
    \begin{equation}
        -\int_{\Omega}\dfrac{A'(\chi)(\hat{\chi})}{A(\chi)^{2}}{E}_0\varepsilon(u):\varepsilon(\hat{u})\ dx+\int_{\Omega}\dfrac{1}{A(\chi)}{E}_0\varepsilon(u'):\varepsilon(\hat{u})\ dx=0,\nonumber
    \end{equation}
    \begin{align*}
        2\int_{\Omega}\dfrac{[A'(\chi)(\hat{\chi})]^2}{A(\chi)^{3}}{E}_0\varepsilon(u):\varepsilon(\hat{u})\ dx-2\int_{\Omega}\dfrac{A'(\chi)(\hat{\chi})}{A(\chi)^2}{E}_0\varepsilon(u'):\varepsilon(\hat{u})\ dx
        +\int_{\Omega}\dfrac{1}{A(\chi)}{E}_0\varepsilon(u''):\varepsilon(\hat{u})\ dx=0
    \end{align*}
   for any $\hat{u}\in H_D^1(\Omega)^d$.
    Hence, setting $\hat{u} = u'$ in the first equation yields
    \begin{equation}\label{state_tangent}
        \int_{\Omega}\dfrac{1}{A(\chi)}{E}_0\varepsilon(u'):\varepsilon(u')\ dx=\int_{\Omega}\dfrac{A'(\chi)(\hat{\chi})}{A(\chi)^{2}}{E}_0\varepsilon(u):\varepsilon(u')\ dx,
    \end{equation} 
    while setting $\hat{u} = u$ in the second equation gives
    \begin{align*}
        \int_{\Omega}\dfrac{1}{A(\chi)}{E}_0\varepsilon(u''):\varepsilon(\hat{u})\ dx
        =-2\int_{\Omega}\dfrac{[A'(\chi)(\hat{\chi})]^2}{A(\chi)^{3}}{E}_0\varepsilon(u):\varepsilon(u)\ dx+2\int_{\Omega}\dfrac{A'(\chi)(\hat{\chi})}{A(\chi)^2}{E}_0\varepsilon(u'):\varepsilon(u)\ dx.
    \end{align*}
    By applying the Cauchy-Schwarz inequality in \eqref{state_tangent}, we obtain
    \begin{equation}
        \int_{\Omega}\dfrac{1}{A(\chi)}{E}_0\varepsilon(u'):\varepsilon(u')\ dx\leq \int_{\Omega}\dfrac{[A'(\chi)(\hat{\chi})]^2}{A(\chi)^{3}}{E}_0\varepsilon(u):\varepsilon(u)\ dx.
    \end{equation}
    Therefore
    \begin{align*}
      \dfrac{\partial^2 g_1(\chi,u,v)}{\partial \chi^2}(\hat{\chi})(\hat{\chi})=-2\int_{\Omega}\dfrac{[A'(\chi)(\hat{\chi})]^2}{A(\chi)^{3}}{E}_0\varepsilon(u):\varepsilon(u)\ dx+2\int_{\Omega}\dfrac{1}{A(\chi)}{E}_0\varepsilon(u'):\varepsilon(u')\ dx\leq 0.
    \end{align*}
This finishes the proof.
\end{proof}

\end{document}
